\section{Related Work and Research Gaps}
\label{sec:related}

In this section, we review the literature on vision-based waste recognition, on waste datasets from India and South Asia, and on hyperparameter optimization. The research gaps are stated at the end of the section.

\subsection{Review Method}
\label{subsec:review-method}

The literature search was carried out in October 2026. Crossref, arXiv, Europe PMC, PubMed Central, OpenAlex and publisher websites were searched with keywords related to waste reviews, datasets and detectors, to India and South Asia, and to metaheuristic and hyperparameter optimization. We also followed the citations of six seed papers~\cite{majchrowska2022deep,yang2016classification,kaya2023classification,sayed2024intelligent,aguerchi2024cnn,alshammari2025smart}, all of which are discussed below. Vision studies on waste or litter and optimization studies on convolutional neural networks (CNNs) were included. Peer-reviewed papers published between 2016 and 2026 were preferred, although the classic method papers are older, and a few preprints, technical reports and workshop papers were retained where no peer-reviewed source covered the topic; their status is mentioned where they are cited. Fifteen studies were examined in detail and are compared in Table~\ref{tab:lit-matrix}, while the remaining works are cited for their stated scope. Since the review is targeted and not systematic, the statements made below about ``the literature'' apply only to the works we reviewed.

\subsection{Vision-Based Waste Recognition: From Classification to Detection}
\label{subsec:rw-waste}

A large share of waste-sorting research classifies images that contain a single object. The TrashNet dataset was released by Yang and Thung~\cite{yang2016classification} as part of a Stanford course report. Each of its 2,527 public images shows one item placed on a white posterboard and belongs to one of six classes. In that report, a support vector machine (SVM) with SIFT features achieved a test accuracy of 63\%, whereas a CNN trained from scratch reached only 22\%. The authors attributed the poor CNN result to unsuitable hyperparameters and left a more thorough hyperparameter search, as well as multi-object detection, for future work. Kaya~\cite{kaya2023classification} truncated pretrained Xception and InceptionResNetV2 networks and obtained 89.72\% test accuracy with one fixed training setting, noting the plain backgrounds of TrashNet as a limitation. Sayed et al.~\cite{sayed2024intelligent} used an InceptionV3 tuned by a metaheuristic, which is discussed in Section~\ref{subsec:rw-hpo}. In the smart bin of Alshammari et al.~\cite{alshammari2025smart}, items are routed by a MobileNetV3-based classifier running on a Raspberry Pi. TrashBox~\cite{kumsetty2022trashbox} carries the word ``detection'' in its title but is presented as a classification dataset of 17,785 images in seven classes. A single label per image is adequate when items are dropped one at a time into a chute. It is not adequate for a heap, a conveyor or a street.

The review papers come to a similar conclusion. Abdallah et al.~\cite{abdallah2020artificial} surveyed 85 studies on artificial intelligence in solid waste management published between 2004 and 2019. Lu and Chen~\cite{lu2022computer} observed that many vision studies on waste sorting are preliminary, are conducted in simplified environments and use artificially collected data; they recommended real-world complexity and open datasets. Fang et al.~\cite{fang2023artificial} reported that most of the vision studies covered by them classify whole images by material, that detection studies are fewer and less accurate, and that laboratory pictures differ from real waste. The accuracies collected in their review range from 72.8\% to 99.95\%, but classification and detection results are mixed together, so the methods cannot be ranked from these numbers.

Datasets for detection and segmentation bring the problem closer to cluttered scenes. TACO~\cite{proenca2020taco}, a crowd-sourced dataset described in a preprint, contained 1,500 photographs of litter in the wild with 4,784 annotations when it was released. ZeroWaste~\cite{bashkirova2022zerowaste} was recorded on the sorting line of a recycling facility and contains deformable and translucent objects in heavy clutter, and WaRP~\cite{yudin2024hierarchical} addresses hierarchical waste detection in images from a recycling plant. UAVVaste~\cite{kraft2021autonomous} consists of 772 low-altitude drone images with 3,716 annotations. Its authors pointed out that the large, near-ground objects of TACO are not suitable for aerial detection. PlastOPol~\cite{cordova2022litter} offers 2,418 crowd-sourced images of urban areas, beaches and forests with 5,300 annotations of a single ``litter'' class. The TrashCan technical report~\cite{hong2020trashcan} provides boxes and masks for 7,212 deep-sea ROV images. FloW~\cite{cheng2021flow} covers floating waste in inland waters and is partly supported by radar data, while MJU-Waste~\cite{wang2020multilevel} contains 2,475 RGB-D laboratory images of people holding waste, labeled with one ``waste'' class. Aerial, underwater, conveyor, hand-held and street settings thus all have their own datasets. Many of these have only one generic class, and their class definitions do not match one another.

The pipeline closest to ours is that of Majchrowska et al.~\cite{majchrowska2022deep}, from which this work started. They merged public datasets into a one-class ``detect-waste'' set and a seven-category ``classify-waste'' set based on the sorting rules of Gda\'{n}sk. Because direct multi-class detection gave poor results, they localized litter with EfficientDet-D2 and classified each crop with EfficientNet-B2. Their hyperparameters were set manually, and one result per source dataset was reported on a held-out 20\% described as a validation and testing set. Fulton et al.~\cite{fulton2019robotic} compared YOLOv2, Tiny-YOLO, Faster R-CNN and SSD for underwater litter on desktop and embedded hardware. Melinte et al.~\cite{melinte2020deep} annotated TrashNet with boxes and reported a precision of 97.63\% with SSD. Liao and Juang~\cite{liao2022realtime} ran YOLOv4-Tiny-3l at 22 frames per second on an Xavier NX to monitor coastal litter from a drone. In a workshop paper, Bartolo et al.~\cite{bartolo2025learning} trained five detectors with privileged information on one drone dataset and tested them on two other datasets. S\'{a}nchez-Ferrer et al.~\cite{sanchezferrer2023experimental} studied Mask R-CNN for marine debris when data are scarce and added synthetic scenes. In all these studies, the hyperparameters are treated merely as part of the experimental setup (Table~\ref{tab:lit-matrix}).

The detector used in SwarmSort belongs to the single-stage family that began with YOLO~\cite{redmon2016you}. This family is frequently chosen in waste studies where speed is important~\cite{kumar2021novel,kraft2021autonomous,liao2022realtime,cordova2022litter}. In the PlastOPol study, for example, YOLOv5x was the most accurate of seven detectors and YOLOv5s ran at up to 5.19 frames per second on a phone. We use YOLOv8n~\cite{jocher2023ultralytics} initialized with COCO-pretrained weights~\cite{lin2014microsoft} (Section~\ref{subsec:detector}).

\subsection{Waste Data From India and South Asia}
\label{subsec:rw-india}

Image data on waste from India are scattered, and most of them were not created for detection. The earliest work we came across is SpotGarbage~\cite{mittal2016spotgarbage} from IIT Ropar, a smartphone application that detects and coarsely segments garbage in geo-tagged photographs so that citizens can report it. Kumar et al.~\cite{kumar2021novel} motivated their study of YOLOv3 and YOLOv3-tiny with Indian waste statistics. They collected 6,437 images in six classes and reported 94.99\% and 51.95\% mAP on a split that served for both validation and testing. Their images, however, are mostly clean and show one or a few items each, and the place where they were taken is not mentioned. AquaVision~\cite{panwar2020aquavision}, written mainly by Indian authors, detects waste in water bodies using a subset derived from TACO. TrashBox~\cite{kumsetty2022trashbox} from NIT Karnataka is a classification dataset, as mentioned earlier.

DWSD~\cite{ali2025dwsd} comes closest to an Indian street benchmark. It contains 784 images taken with phone cameras in streets, parks, lawns and garbage areas around Kolkata (Jadavpur, Garia) and Berhampore in West Bengal, with pixel-level masks for 14 waste classes. The baselines in the paper report only pixel accuracy (89.39\% for DeepLabv3+), and no detection results are given. We also noticed that the paper is not consistent about the name of the dataset and the number of images. From the rest of South Asia, UGV-NBWASTE~\cite{isalm2025ugvnbwaste} provides 3,600 images of eight types of non-biodegradable waste from Barisal, Bangladesh. It has oriented boxes, but only about 1.1 objects per image and no baseline models. The pLitterStreet preprint~\cite{mandhati2024plitterstreet} maps plastic litter at street level in Thailand, Vietnam and Sri Lanka using cameras mounted on vehicles. Fang et al.~\cite{fang2023artificial} also mention the scarcity of data in developing countries as one of the limitations of the field.

\subsection{Hyperparameter Optimization and Swarm Intelligence}
\label{subsec:rw-hpo}

Particle swarm optimization (PSO) was proposed by Kennedy and Eberhart~\cite{kennedy1995particle}. A population of candidate solutions, called particles, moves through the search space, and the velocity of each particle is a combination of its previous velocity and the attraction toward its own best position and the best position of the swarm. Shi and Eberhart~\cite{shi1998modified} later added the inertia weight, which scales the previous velocity. Clerc and Kennedy~\cite{clerc2002particle} analyzed the particle trajectories and derived coefficients that control convergence, which led to constriction-factor PSO. PSO does not require gradients and evaluates a fixed number of candidates in every round. This makes it convenient when each evaluation is expensive, as in the case of training a network. The settings used in our work are given in Section~\ref{subsec:pso}.

Random search~\cite{bergstra2012random} is the obvious reference for any guided search method. It samples the settings independently and can be parallelized without restriction. If an optimizer cannot do better than random search with the same number of evaluations, there is no proof that its guidance is useful. Multi-fidelity methods take a different route and reduce the cost of each evaluation. Hyperband~\cite{li2018hyperband}, for example, accelerates random search by allocating more resources to promising configurations and stopping the weak ones early. SwarmSort uses a simpler approach with one fixed low fidelity, namely a short proxy training (Section~\ref{subsec:proxy}).

Metaheuristics have been applied to tune CNNs in several domains, but we found the supporting evidence to be thin. Aguerchi et al.~\cite{aguerchi2024cnn} used PSO to select three architecture hyperparameters (number of filters, kernel size and stride) of a small CNN for mammogram masses, while the training hyperparameters were kept fixed. Their swarm had 4 particles and ran for 3 iterations, which amounts to roughly 12 network evaluations. No baseline (default setting, random search or another optimizer) was included and only a single result was reported, even though the authors themselves remark that a metaheuristic requires repeated runs and statistical analysis. The most complete comparison in the waste domain that we found is by Sayed et al.~\cite{sayed2024intelligent}. A multi-objective beluga whale optimizer (MBWO) tunes the learning rate, batch size and dropout period of InceptionV3 on TrashNet with a population of 30, 20 iterations and 15 epochs per evaluation. Four other metaheuristics are run with the same population, iterations, bounds and fitness function, and an untuned model serves as reference (90.76\% accuracy). A Wilcoxon rank-sum test on the best search fitness over independent runs shows that MBWO is significantly better than three of the competitors but not the fourth (EO). On the other hand, the number of runs is not stated, random search is not included, the test metrics are single values without intervals, and the headline accuracy is 97.75\% in the abstract but 97.57\% in the results table. Mao et al.~\cite{mao2021recycling} applied a genetic algorithm to a CNN that classifies recycling waste.

\subsection{Comparative Analysis and Research Gaps}
\label{subsec:gaps}

Table~\ref{tab:lit-matrix} compares the 15 studies that we examined in detail with respect to the points that determine whether a tuning claim is reliable. Only two of them tune hyperparameters with a metaheuristic. Both are classifiers, and only one compares tuning methods. None of the studies reports a confidence interval on a test metric, and only Bartolo et al.~\cite{bartolo2025learning} test on public datasets that were not seen during training. Based on these observations and on Sections~\ref{subsec:rw-waste}--\ref{subsec:rw-hpo}, we identify four gaps. For each gap, the evidence that goes against it is also mentioned.

\begin{table*}[!t]
\centering
\caption{Comparison of Existing Studies With SwarmSort}
\label{tab:lit-matrix}
\footnotesize
\setlength{\tabcolsep}{3pt}
\renewcommand{\arraystretch}{1.1}
\begin{tabular}{@{}>{\raggedright\arraybackslash}p{2.75cm}>{\raggedright\arraybackslash}p{1.45cm}>{\raggedright\arraybackslash}p{1.75cm}>{\raggedright\arraybackslash}p{2.2cm}>{\raggedright\arraybackslash}p{2.75cm}>{\raggedright\arraybackslash}p{2.45cm}>{\raggedright\arraybackslash}p{1.95cm}>{\raggedright\arraybackslash}p{1.3cm}@{}}
\toprule
Study & Task & Scenes & Hyperparameters set by & Tuning baselines & Runs / CIs & Cross-domain test & India data \\
\midrule
Yang \& Thung~\cite{yang2016classification} & Cls & Studio & Manual & --- & 1 / no & No & No \\
Kaya~\cite{kaya2023classification} & Cls & TrashNet & Fixed & --- & 1 / no & No & No \\
Sayed et al.~\cite{sayed2024intelligent} & Cls & TrashNet & MBWO & 4 metaheur.\ (equal), untuned; no RS & n.r.\ + Wilcoxon (fitness) / no & No & No \\
Aguerchi et al.~\cite{aguerchi2024cnn} & Cls & Clinical & PSO (architecture) & None & 1 / no & No & --- \\
Majchrowska et al.~\cite{majchrowska2022deep} & Det + Cls & Real-world & Manual & None & 1 / no & No & No \\
Hong et al.~\cite{hong2020trashcan} & Det / Seg & Underwater & Fixed & None & 1 / no & No & No \\
Kraft et al.~\cite{kraft2021autonomous} & Det & Aerial & Not described & None & 1 / no & No & No \\
C\'{o}rdova et al.~\cite{cordova2022litter} & Det & Real-world & Published defaults & None & No SDs / no & Per dataset & No \\
Wang et al.~\cite{wang2020multilevel} & Seg & Lab + TACO & Fixed; CRF grid & None & 1 / no & No & No \\
Fulton et al.~\cite{fulton2019robotic} & Det & Underwater & Community & None & 1 / no & Shifted footage & No \\
Melinte et al.~\cite{melinte2020deep} & Det & TrashNet & Manual + LR test & None & 1 / no & Lab-made set & No \\
Kumar et al.~\cite{kumar2021novel} & Det & Mostly clean & Fixed & None & 1 / test = val & No & Unstated \\
Liao \& Juang~\cite{liao2022realtime} & Det & Aerial & ``Tuned'' (n.r.) & None & No & Qualitative & No \\
Bartolo et al.~\cite{bartolo2025learning} & Det & Aerial & Manual & None & 1 / no & Yes & No \\
Ali et al.~\cite{ali2025dwsd} & Seg & Real-world & n.r. & --- & 1 / no & No & Yes \\
\midrule
\textbf{SwarmSort (ours)} & \textbf{Det} & \textbf{Studio + real-world} & \textbf{PSO vs RS vs defaults} & \textbf{Equal: 32 proxy trainings each} & \textbf{1 seed / bootstrap 95\% CIs} & \textbf{Yes (3 domains)} & \textbf{DWSD (test)} \\
\bottomrule
\end{tabular}

\vspace{3pt}
\begin{minipage}{\textwidth}
\scriptsize
Datasets: TrashNet~\cite{yang2016classification} (single-object studio images), TrashCan~\cite{hong2020trashcan}, UAVVaste~\cite{kraft2021autonomous}, PlastOPol~\cite{cordova2022litter}, MJU-Waste~\cite{wang2020multilevel} and DWSD~\cite{ali2025dwsd} (a dataset paper). Cls = classification; Det = detection (one class in~\cite{majchrowska2022deep,cordova2022litter}); Seg = segmentation; MBWO = multi-objective beluga whale optimizer; metaheur.\ (equal) = other metaheuristics with the same population size and iterations; RS = random search; CRF grid = grid search for the conditional random field only; LR test = learning-rate range test; SD = standard deviation; CI = confidence interval; n.r.\ = not reported; --- = not applicable. ``1 / no'' = one reported run per configuration and no confidence intervals. ``Fixed'' = one setting used without a described selection procedure. ``test = val'' = the split used for validation also served as the test split. Per dataset = two datasets, each trained and tested separately; Qualitative = new sites described only qualitatively; Unstated = Indian authors, image location not stated. SwarmSort: 6 classes in \genone and 7 in \gentwo; one training seed per arm; paired-bootstrap CIs; studio, real-world and Indian test sets. The review of Fang et al.~\cite{fang2023artificial} and the UGV-NBWASTE dataset~\cite{isalm2025ugvnbwaste}, which reports no model, are not listed.
\end{minipage}
\end{table*}

\textbf{Gap 1 (task and scene).} The classifiers discussed in Section~\ref{subsec:rw-waste}~\cite{yang2016classification,kaya2023classification,sayed2024intelligent,alshammari2025smart,kumsetty2022trashbox} give one label per image, and those based on TrashNet see one object on a plain background. Fang et al.~\cite{fang2023artificial} found whole-image classification to be the most common approach. Even among the detection studies, some depend on data that have nearly one object per image~\cite{melinte2020deep,kumar2021novel,isalm2025ugvnbwaste}. At the same time, multi-object waste detection is clearly an active line of work, mostly in computer-vision, robotics and sensing venues, as the datasets and the pipeline in Section~\ref{subsec:rw-waste} show~\cite{proenca2020taco,bashkirova2022zerowaste,cheng2021flow,yudin2024hierarchical,kraft2021autonomous,cordova2022litter,hong2020trashcan,majchrowska2022deep}. As our review is not systematic, we do not attempt to quantify how much smaller this line is. The gap can be stated as follows: \emph{in the reviewed literature, waste recognition for sorting is dominated by single-label classification of one object on a plain background; multi-object detection in cluttered scenes is a smaller strand, and several detection studies still rely on single-object studio images.}

\textbf{Gap 2 (optimization).} In waste vision, metaheuristic tuning has been applied to classifiers~\cite{sayed2024intelligent,mao2021recycling}, and PSO tuning of CNNs is found in other domains~\cite{aguerchi2024cnn}. None of the waste detection or segmentation studies in Table~\ref{tab:lit-matrix} describes a search over training hyperparameters. The only search that is described concerns post-processing parameters~\cite{wang2020multilevel}, and Fang et al.~\cite{fang2023artificial} mention PSO and genetic algorithms only in connection with routing and prediction. The gap is therefore limited to the combination of metaheuristic tuning with waste detectors. Two detection papers mention tuning without describing it~\cite{liao2022realtime,bartolo2025learning}, so a search in those cases cannot be excluded, and for the works that were not examined in detail the training setup could not be verified. The gap is: \emph{in the reviewed literature, no waste detector had its training hyperparameters tuned by a swarm optimizer; detection studies used default or hand-set values, or a tuning step that is not described.}

\textbf{Gap 3 (evidence).} When tuning is evaluated, the baselines are weak. Aguerchi et al.~\cite{aguerchi2024cnn} have none, and Sayed et al.~\cite{sayed2024intelligent} compare only with other metaheuristics; in both cases the test metrics are single numbers (Section~\ref{subsec:rw-hpo}). Among the detection and segmentation papers in Table~\ref{tab:lit-matrix}, none reports run-to-run variation, confidence intervals or significance tests for its accuracy, and Bartolo et al.~\cite{bartolo2025learning} state that statistical analysis was not applicable. It must be acknowledged that Sayed et al. do compare optimizers with equal population and iterations, include an untuned reference and apply a significance test to the search fitness. We therefore do not claim that optimizers are never compared or that significance is never tested. The gap is: \emph{metaheuristic tuning in this area is compared against other metaheuristics or against nothing; we found no comparison with random search~\cite{bergstra2012random} at an equal evaluation budget, and test results are reported as single numbers without confidence intervals.} SwarmSort reduces this gap but does not close it completely. Since we use one training seed per arm and one search seed per method, our intervals account only for the sampling of the test set (Section~\ref{subsec:limitations}).

\textbf{Gap 4 (generalization).} The TrashNet-based classifiers are never evaluated outside the studio, two reviews point to simplified, laboratory-like data~\cite{lu2022computer,fang2023artificial}, and accuracy depends strongly on the dataset, as the PlastOPol example in Section~\ref{sec:intro} shows~\cite{cordova2022litter}. Most detection studies test within a single dataset or describe the degradation at new sites only qualitatively (Table~\ref{tab:lit-matrix}), and the Indian work in Section~\ref{subsec:rw-india} does not provide a detection benchmark. Cross-domain tests are not entirely absent. Bartolo et al.~\cite{bartolo2025learning} test on datasets not seen in training, Fulton et al.~\cite{fulton2019robotic} on footage that differs in time and place, and S\'{a}nchez-Ferrer et al.~\cite{sanchezferrer2023experimental} add synthetic scenes to scarce real data. South Asian detection data are available in UGV-NBWASTE~\cite{isalm2025ugvnbwaste}, although without baselines, and in the pLitterStreet preprint~\cite{mandhati2024plitterstreet}. The gap is: \emph{studio-to-real-world and cross-dataset evaluation is rare and recent, and we found no Indian real-world waste test set with published object-detection baselines.} SwarmSort reports studio, real-world and DWSD results separately and compares studio-only training with mixed training on the same test images (Section~\ref{subsec:res-domain}). As already mentioned, the DWSD boxes are obtained from mask regions that merge touching items and the class mapping is inferred from the masks, so the DWSD results are affected by a labeling mismatch as well.

\textbf{Observation O5 (reporting of positive results).} The applied papers that we read present their methods as improvements. Sayed et al.~\cite{sayed2024intelligent}, for example, describe their optimizer as the best although its advantage over EO is not significant. Some negative findings are reported in part: a CNN that performs far below an SVM~\cite{yang2016classification}, direct multi-class detection being abandoned because of poor results~\cite{majchrowska2022deep}, synthetic data that help with data scarcity only to a limited extent~\cite{sanchezferrer2023experimental}, and lower performance on shifted data~\cite{fulton2019robotic,liao2022realtime}. We did not, however, find a paper reporting that a tuning method failed to beat a simple baseline. Since the absence of such papers in a targeted sample proves little, we list this as an observation and not as a gap. In SwarmSort the null and negative results are reported and analyzed (Sections~\ref{sec:results} and~\ref{sec:discussion}). Table~\ref{tab:gap-map} shows how each gap is addressed in this work.

\begin{table}[!t]
\centering
\caption{Research Gaps and How SwarmSort Addresses Them}
\label{tab:gap-map}
\footnotesize
\setlength{\tabcolsep}{3pt}
\renewcommand{\arraystretch}{1.1}
\begin{tabular}{@{}>{\raggedright\arraybackslash}p{1.45cm}>{\raggedright\arraybackslash}p{4.85cm}>{\raggedright\arraybackslash}p{1.95cm}@{}}
\toprule
Gap & How SwarmSort addresses it & Where \\
\midrule
G1: task and scene & Multi-object detection with YOLOv8n (6 or 7 classes) on studio and real-world scenes & Sec.~\ref{subsec:datasets}, \ref{subsec:detector} \\
G2: optimization & Our PSO implementation tunes 6 training hyperparameters of a waste detector & Sec.~\ref{subsec:space}, \ref{subsec:pso} \\
G3: evidence & Defaults, random search and PSO at 32 proxy trainings each; proxy check with a threshold fixed in advance (\genone only); test split used once; paired-bootstrap 95\% CIs & Sec.~\ref{subsec:design}, \ref{subsec:proxy}, \ref{subsec:eval}, \ref{subsec:res-gen1}, \ref{subsec:res-gen2} \\
G4: generalization & Real-world training data in \gentwo; studio, real-world and India (DWSD, test only, boxes from masks) results reported separately; both generations on the same test images & Sec.~\ref{subsec:datasets}, \ref{subsec:res-domain} \\
O5: reporting & Null and negative results reported and explained, with limitations & Sec.~\ref{sec:results}, \ref{sec:discussion} \\
\bottomrule
\end{tabular}
\end{table}
